% Ch4 WS2 -- solubility rules, precipitation, ionic equations. Blocks 2-3.
\documentclass{apchem-worksheet}

\apcunitword{Chapter}
\apcunit{4}
\apcunittitle{Types of Chemical Reactions and Solution Stoichiometry}
\apcdoctitle{Worksheet 2 \textbullet\ Precipitation and Net Ionic Equations}
\apcsubtitle{Zumdahl \S4.5--4.7 \textbullet\ cite the rule number for every solubility call}

\begin{document}
\wsheader[For every solubility decision, name the rule that decided it.
For every net ionic equation, list the spectator ions separately.]

\problem[]{Soluble or insoluble? Cite the rule number.}
\begin{parts}
  \item \ce{Na2CO3}: \answerblank[5cm]{soluble --- rule 2}
  \item \ce{AgBr}: \answerblank[5.9cm]{insoluble --- rule 3 exception}
  \item \ce{Ba(NO3)2}: \answerblank[5cm]{soluble --- rule 1}
  \item \ce{PbSO4}: \answerblank[5.9cm]{insoluble --- rule 4 exception}
  \item \ce{Mg(OH)2}: \answerblank[5cm]{insoluble --- rule 5}
  \item \ce{(NH4)3PO4}: \answerblank[5.9cm]{soluble --- rule 2 beats rule 6}
\end{parts}

\problem[]{For each pair of solutions, predict the products by ion
interchange, decide whether a precipitate forms, and write the balanced
formula equation (or ``no reaction'').}
\begin{parts}
  \item \ce{Pb(NO3)2(aq) + 2 KI(aq)}
        \answerlines[2]{\ce{Pb(NO3)2(aq) + 2 KI(aq) -> PbI2(s) +
        2 KNO3(aq)} --- \ce{PbI2} insoluble by rule 3.}
  \item \ce{NaNO3(aq) + KCl(aq)}
        \answerlines[2]{No reaction --- both \ce{NaCl} and \ce{KNO3} are
        soluble (rules 1--3), so no solid forms.}
  \item \ce{BaCl2(aq) + Na2SO4(aq)}
        \answerlines[2]{\ce{BaCl2(aq) + Na2SO4(aq) -> BaSO4(s) +
        2 NaCl(aq)} --- \ce{BaSO4} insoluble by rule 4.}
\end{parts}

\problem[]{For \ce{AgNO3(aq) + NaCl(aq)}:}
\begin{parts}
  \item Formula equation: \answerlines[1]{\ce{AgNO3(aq) + NaCl(aq) ->
        AgCl(s) + NaNO3(aq)}}
  \item Complete ionic equation: \answerlines[2]{\ce{Ag+(aq) + NO3^-(aq) +
        Na+(aq) + Cl^-(aq) -> AgCl(s) + Na+(aq) + NO3^-(aq)}}
  \item Spectator ions: \answerblank[4.5cm]{\ce{Na+} and \ce{NO3-}}
  \item Net ionic equation: \answerblank[6.6cm]{\ce{Ag+(aq) + Cl^-(aq) ->
        AgCl(s)}}
\end{parts}

\problem[]{Write the net ionic equation for each:}
\begin{parts}
  \item \ce{Na2CO3(aq) + CaCl2(aq)}
        \answerblank[7cm]{\ce{Ca^2+ + CO3^2- -> CaCO3(s)}}
  \item \ce{K2S(aq) + Cu(NO3)2(aq)}
        \answerblank[7cm]{\ce{Cu^2+ + S^2- -> CuS(s)}}
  \item \ce{(NH4)2SO4(aq) + Ba(OH)2(aq)}
        \answerlines[2]{\ce{Ba^2+ + SO4^2- -> BaSO4(s)} (and, if the
        ammonium/hydroxide reaction is included,
        \ce{NH4+ + OH^- -> NH3 + H2O}).}
\end{parts}

\problem[]{Three different reactions all give the net ionic equation
\ce{Ba^2+(aq) + SO4^2-(aq) -> BaSO4(s)}. Write formula equations for two
of them, using different spectator ions.
\answerlines[3]{e.g.\ \ce{BaCl2(aq) + Na2SO4(aq) -> BaSO4(s) + 2 NaCl(aq)}
and \ce{Ba(NO3)2(aq) + K2SO4(aq) -> BaSO4(s) + 2 KNO3(aq)}.}}

\problem[]{A student writes this net ionic equation for mixing
\ce{HC2H3O2(aq)} with \ce{NaOH(aq)}:
\ce{H+ + OH^- -> H2O}. Identify the error and give the correct equation.
\answerlines[3]{Acetic acid is a weak electrolyte, so it must be written
intact rather than as \ce{H+}. Correct:
\ce{HC2H3O2(aq) + OH^-(aq) -> C2H3O2^-(aq) + H2O(l)}.}}

\problem[]{Precipitation stoichiometry. \SI{40.0}{\milli\liter} of
\SI{0.250}{\Molar} \ce{Pb(NO3)2} is mixed with excess \ce{KI}.}
\begin{parts}
  \item Moles of \ce{Pb^2+}: \answerblank[3cm]{\SI{0.0100}{\mole}}
  \item Mass of \ce{PbI2} formed ($M = 461.0$): \workspace[2cm]
        \keyonly{\begin{apcanswer}1:1 ratio $\to 0.0100 \times 461.0 =
        \SI{4.61}{\gram}$\end{apcanswer}}
  \item Moles of \ce{KI} required at minimum:
        \answerblank[3cm]{\SI{0.0200}{\mole}}
\end{parts}

\begin{spiralstrip}{Chapter 4 \textbullet\ block 1}
  \item $[\ce{Cl-}]$ in \SI{0.10}{\Molar} \ce{AlCl3}:
        \answerblank[2.6cm]{\SI{0.30}{\Molar}}
  \item Strong, weak, or non: \ce{NH3}? \answerblank[2.6cm]{weak}
  \item Dilute \SI{25.0}{\milli\liter} of \SI{2.00}{\Molar} to
        \SI{100.}{\milli\liter}: \answerblank[2.6cm]{\SI{0.500}{\Molar}}
  \item Do you split water in an ionic equation?
        \answerblank[2cm]{no}
  \item Which ion makes almost everything soluble?
        \answerblank[3.8cm]{nitrate (or \ce{NH4+})}
\end{spiralstrip}

\end{document}
